\subsection{SM4 对称加密算法}

SM4 是我国商用密码标准中的分组对称加密算法，于 2016 年作为国家标准 GB/T 32907-2016 发布\footnote{GB/T 32907-2016《信息安全技术\ SM4分组密码算法》，国家标准化管理委员会，2016年8月29日发布，2017年3月1日实施。}。SM4 的分组长度为 128 bit，密钥长度为 128 bit，采用 32 轮非平衡 Feistel 结构，每轮使用相同的轮函数但以不同的轮密钥进行运算。SM4 的加密和解密使用相同的算法结构，区别仅在于轮密钥的使用顺序相反。对称加密算法的特点是加密和解密使用相同密钥，具有计算效率高、适合批量数据加密等优点，常用于通信数据保护、数据库字段加密和文件加密等场景。

设$K$表示对称加密密钥，$M$表示明文消息，$C$表示密文消息，则 SM4 加密与解密过程可以抽象表示为：

\begin{equation}
C = \operatorname{Enc}_{K}(M)
\end{equation}

\begin{equation}
M = \operatorname{Dec}_{K}(C)
\end{equation}

SM4 的轮函数由以下组件构成：输入分组的 4 个 32 bit 字中，第一个字与其余三个字经轮函数$F$处理后异或，再与轮密钥结合，生成下一轮的输入。轮函数$F$中包含一个 8 bit 输入/8 bit 输出的 S 盒（非线性置换）和一个线性变换$L$（由循环移位和异或组成）。密钥扩展算法使用与加密算法类似的结构，将 128 bit 主密钥扩展为 32 个 32 bit 轮密钥。

在实际系统中，SM4 通常需要结合工作模式使用，以处理长度大于一个分组的数据并增强安全性。常见的工作模式包括：

\begin{itemize}[leftmargin=2em]
\item \textbf{ECB（Electronic Codebook）模式}：每个分组独立加密，相同明文分组产生相同密文，不适合加密结构化数据。
\item \textbf{CBC（Cipher Block Chaining）模式}：每个明文分组在加密前先与前一个密文分组异或，引入初始化向量$IV$使相同明文在不同加密过程中得到不同密文，安全性优于 ECB 模式。
\item \textbf{CTR（Counter）模式}：将递增计数器作为加密输入，输出与明文异或得到密文，支持并行加密。
\item \textbf{GCM（Galois/Counter Mode）模式}：在 CTR 模式基础上增加认证标签，同时提供加密和完整性保护。
\end{itemize}

以 CBC 模式为例，加密时需要引入初始化向量$IV$，使相同明文在不同加密过程中得到不同密文，从而降低明文模式泄露风险。其抽象形式如下：

\begin{equation}
C = \operatorname{SM4\text{-}CBC\text{-}Enc}_{K}(IV, M)
\end{equation}

具体而言，CBC 模式的加密过程为：$C_0 = IV$，$C_i = \operatorname{Enc}_K(M_i \oplus C_{i-1})$，其中$M_i$和$C_i$分别表示第$i$个明文分组和密文分组。解密过程为：$M_i = \operatorname{Dec}_K(C_i) \oplus C_{i-1}$。由于每个分组的加密都依赖于前一个密文分组，攻击者无法仅通过替换某个密文分组来篡改对应的明文而不影响后续分组的解密结果。

% !TEX root = ../../main.tex
% SM4-CBC 加密链图
\begin{figure}[H]
\centering
\begin{tikzpicture}[
  node distance = 0.8cm and 1.6cm,
  blk/.style  = {rectangle, draw, rounded corners=2pt, minimum width=1.4cm, minimum height=0.55cm, font=\small, fill=blue!10, draw=blue!50!black},
  enc/.style  = {rectangle, draw, minimum width=1.4cm, minimum height=0.6cm, font=\small\bfseries, fill=orange!20, draw=orange!70!black},
  xop/.style  = {circle, draw, minimum size=0.45cm, font=\small\bfseries, fill=gray!15},
  ct/.style   = {rectangle, draw, rounded corners=2pt, minimum width=1.4cm, minimum height=0.55cm, font=\small, fill=green!12, draw=green!50!black},
  arr/.style  = {-Stealth, thick},
  lbl/.style  = {font=\scriptsize\itshape},
  keylbl/.style = {font=\scriptsize\bfseries, fill=white, inner sep=1pt},
]

% IV
\node[blk] (IV)  {$IV$};
% M1
\node[blk, right=1.6cm of IV]  (M1) {$M_1$};
% M2
\node[blk, right=2.0cm of M1] (M2) {$M_2$};
% M3
\node[blk, right=2.0cm of M2] (M3) {$M_3$};

% XOR 1
\node[xop, below=0.55cm of M1] (x1) {$\oplus$};
% XOR 2
\node[xop, below=0.55cm of M2] (x2) {$\oplus$};
% XOR 3
\node[xop, below=0.55cm of M3] (x3) {$\oplus$};

% Enc 1
\node[enc, below=0.55cm of x1] (e1) {SM4-Enc};
% Enc 2
\node[enc, below=0.55cm of x2] (e2) {SM4-Enc};
% Enc 3
\node[enc, below=0.55cm of x3] (e3) {SM4-Enc};

% C1 C2 C3
\node[ct, below=0.55cm of e1] (C1) {$C_1$};
\node[ct, below=0.55cm of e2] (C2) {$C_2$};
\node[ct, below=0.55cm of e3] (C3) {$C_3$};

% 密钥标注 — 放在 enc 节点右侧，留足间距避免重叠
\node[keylbl, right=0.35cm of e1] {$K$};
\node[keylbl, right=0.35cm of e2] {$K$};
\node[keylbl, right=0.35cm of e3] {$K$};

% 箭头：明文→XOR→Enc→密文
\draw[arr] (M1) -- (x1);
\draw[arr] (M2) -- (x2);
\draw[arr] (M3) -- (x3);
\draw[arr] (x1) -- (e1);
\draw[arr] (x2) -- (e2);
\draw[arr] (x3) -- (e3);
\draw[arr] (e1) -- (C1);
\draw[arr] (e2) -- (C2);
\draw[arr] (e3) -- (C3);

% IV 指向 XOR1
\draw[arr] (IV.south) -- ++(0,-0.18) -| (x1.north);

% C1 → XOR2, C2 → XOR3（从下方绕行）
\draw[arr] (C1.south) -- ++(0,-0.25) -- ++(1.0,0) |- (x2.west);
\draw[arr] (C2.south) -- ++(0,-0.25) -- ++(1.0,0) |- (x3.west);

% 省略号
\node[font=\Large, right=0.5cm of C3] {$\cdots$};

% 顶部标注 — 放在节点上方，留足间距
\node[lbl, above=0.2cm of IV]  {随机初始化向量};
\node[lbl, above=0.2cm of M1]  {明文分组 1};
\node[lbl, above=0.2cm of M2]  {明文分组 2};
\node[lbl, above=0.2cm of M3]  {明文分组 3};

\end{tikzpicture}
\caption{SM4-CBC 工作模式加密链示意图}
\label{fig:sm4_cbc}
\end{figure}


在 FoodShield 系统中，SM4 主要用于订单通信消息的加密存储。用户与骑手之间的消息经过服务器处理后，使用订单级会话密钥或平台侧加密密钥进行加密，再写入数据库。管理员端默认不直接展示消息明文，而是展示消息哈希、时间戳、订单号和 Merkle Root 等审计信息。这样可以降低数据库泄露或管理员越权查看造成的隐私风险。

需要说明的是，本作品原型系统以可运行、可展示和安全机制闭环为主要目标。在实际部署环境中，通信链路还应结合 HTTPS/TLS 保障传输安全，SM4 更多用于服务端存储阶段的消息内容保护。通过"传输层加密 + 存储层加密 + 日志哈希审计"的组合设计，可以从多个层面降低订单通信数据泄露和篡改风险。
